\section{Introduction}\label{sec:introduction}
For $f\in L^2(\R^2)$ we use the Fourier transform
\[
\widehat f(\xi)=\int_{\R^2}e^{-ix\cdot\xi}f(x)\,dx,
\qquad
\|f\|_{H^{1/3}}^2=(2\pi)^{-2}\int_{\R^2}
 (1+|\xi|^2)^{1/3}|\widehat f(\xi)|^2\,d\xi.
\]
The first formula is initially defined on integrable functions and then
extended by continuity to $L^2$. For $a>0$ put
\begin{equation}\label{eq:intro-abel}
 U_af(x,t)=(2\pi)^{-2}\int_{\R^2}
 e^{ix\cdot\xi-it|\xi|^2-a|\xi|^2}\widehat f(\xi)\,d\xi.
\end{equation}
This integral is absolutely convergent by Cauchy--Schwarz. At a Lebesgue
point of $f$, write
\[
 f^*(x)=\lim_{r\downarrow0}\frac1{\pi r^2}
                    \int_{|y-x|<r}f(y)\,dy.
\]
Our main result includes the order of limits in the definition of the
evolution.
\begin{theorem}[Planar endpoint convergence]\label{thm:main}
For every $f\in H^{1/3}(\R^2)$ there is a set $E_f\subset\R^2$ of full
Lebesgue measure such that, for each $x\in E_f$, the finite limit
\[
 u_f(x,t)=\lim_{a\downarrow0}U_af(x,t)
\]
exists for every real $0<t<1$, and
\[
 \lim_{t\downarrow0}u_f(x,t)=f^*(x).
\]
In fact $u_f(x,\cdot)$ extends continuously to $[0,1)$ with value $f^*(x)$
at zero, and the Gaussian limit is uniform over $0<t<1$.
\end{theorem}

\subsection{History and significance}
Carleson \cite{Carleson80} introduced the pointwise convergence problem and
proved its one-dimensional positive result at regularity $s=1/4$.
Dahlberg and Kenig \cite{DahlbergKenig82} showed that this threshold is
necessary. Sj\"olin \cite{Sjolin87} and Vega \cite{Vega88} independently
obtained convergence for $s>1/2$ in higher dimensions. In the plane,
Bourgain and Moyua--Vargas--Vega
\cite{Bourgain95,MVV96} developed the connection with Fourier restriction.
The bilinear approach of Tao and Vargas \cite{TV00} gave $s>15/32$;
Tao's sharp bilinear restriction theorem \cite{Tao03} improved this to
$s>2/5$, and Lee \cite{Lee06} obtained $s>3/8$.
In dimensions $n\ge3$, Bourgain \cite{Bourgain13} obtained the sufficient
range $s>1/2-1/(4n)$.

Bourgain's later counterexample \cite[Proposition~1]{Bourgain16} forces
$s\ge n/[2(n+1)]$ when $n\ge2$, but does not rule out equality.
Du, Guth and Li \cite[Theorems~1.1--1.2]{DGL17} proved planar convergence
and a local $L^3$ maximal estimate for $s>1/3$. Their proof combined
polynomial partitioning with wave packets and refined Strichartz
estimates, using $\ell^2$ decoupling to control packets tangent to the
partitioning surface. Du, Guth, Li and Zhang
\cite[Theorems~1.1 and~4.2]{DGLZ18} then obtained higher-dimensional
convergence for $s>(n+1)/[2(n+2)]$ and developed multilinear refined
Strichartz estimates. Du and Zhang \cite[Theorem~1.1]{DZ19} established
convergence for $s>n/[2(n+1)]$ when $n\ge3$. A recent Bessel-potential
extension still states its positive $p=2$ range with the strict
inequality $s>n/[2(n+1)]$
\cite[Theorems~1.3--1.4]{PanSunTan26}.

The choice of pointwise representative also has a history: Sj\"ogren and
Sj\"olin \cite[Theorem~2]{SjogrenSjolin89} constructed, for $s>1/2$,
representatives continuous on the entire time line for almost every
spatial point, using spatial or space-time averages.

Theorem~\ref{thm:main} resolves positively the planar Sobolev endpoint of
Carleson's convergence problem: it supplies the equality case $s=1/3$,
with the order of limits and the all-real-times representative specified
above. Its maximal estimate has local $L^1$ output. The stronger question
of an endpoint local $L^3$ maximal estimate is not asserted here.

\subsection{Three ingredients of the proof}
The main analytic conclusion is a local maximal estimate at the exact
Sobolev exponent. Write $S(t)f=e^{it\Delta}f$ for the $L^2$ evolution
with Fourier multiplier $e^{-it|\xi|^2}$. For bounded-frequency data,
the pointwise maximal estimate takes the form
\begin{equation}\label{eq:intro-maximal}
 \int_{B(x_0,1)}\sup_{0\le t<1}|S(t)f(x)|\,dx
       \le C\|f\|_{H^{1/3}},
\end{equation}
with a constant independent of the center $x_0$ and the frequency bound.
Theorem~\ref{thm:maximal} proves this estimate. The final section then
constructs the representative in Theorem~\ref{thm:main}; it does not infer
an all-times assertion by intersecting time-dependent exceptional sets.

We work first at unit frequency, with packets of spatial width $L$ and
propagation time $L^2$. A measure satisfying a two-dimensional ball bound
models the spatial area on a measurable time graph. The proof has three
main components.

\paragraph{A transverse power saving.}
Multilinear restriction \cite{BCT06} and the refined fractal restriction
estimate of Du and Zhang \cite[Theorem~1.6]{DZ19} reduce a hypothetical
saturated transverse configuration to weighted incidences between points
and tubes.
Conditional information across scales has additive profiles whose
extremal relations come from three transverse directions. A robust
projection theorem \cite[Theorem~1]{He20} restricts the derivatives of
these profiles to zero
and one. A further argument rules out arbitrarily many alternations;
polynomial partitioning \cite[Theorem~4.1]{GK15} then excludes the
remaining jump. This gives a
fixed power saving, which is propagated to the sparse packet estimate in
Proposition~\ref{prop:sparse}.

\paragraph{A bilinear estimate for smooth packet arrays.}
Theorem~\ref{thm:frame} treats packets with general smooth profiles and
separated frequency groups. Positive masks remove large pieces of a tight
frequency frame, charging them to the lost energy. A quartic potential,
together with the curvature identity for the parabola, controls changes
in the spatial smoothing scale. The resulting estimate is uniform in the
number of scales. This uniformity is needed for the long, nearly flat
frequency configurations that arise in the proof of
Theorem~\ref{thm:distribution}.

\paragraph{Square summation at the endpoint.}
Theorem~\ref{thm:distribution} yields a short-time weak estimate with some
fixed output exponent $p>2$, with no growing frequency loss. To combine
frequencies, we linearize all of them on the same time graph. A binary time
tree separates heavier children from their lighter siblings. Spatial
localization supplies decay in the relative depth of this tree, while
$p>2$ supplies a positive convexity deficit that sums powers of the
lighter-child masses across frequency scales. Together these estimates
retain the Sobolev square norm and prove \eqref{eq:intro-maximal}.

The packet and measure formulations allow arbitrary bounded phase-space
counts and sufficiently smooth compact profiles. In particular, the
bilinear estimate applies to the composite packets created by localization;
no exact evolution equation is imposed on those profiles.

\subsection{Organization and dependencies}
Section~\ref{sec:preliminaries} fixes the packet conventions and the
restriction estimates used with small losses. The transverse and sparse
arguments follow in Sections~\ref{sec:transverse} and~\ref{sec:sparse}.
The independent frame estimate is proved before its applications, in
Section~\ref{sec:frame}. Sections~\ref{sec:flat} and~\ref{sec:geometry}
provide the flat-frequency and geometric reductions, and
Section~\ref{sec:dyadic} proves the uniform dyadic distribution estimate.
Section~\ref{sec:endpoint} performs the endpoint summation and proves
Theorem~\ref{thm:main}. Figure~\ref{fig:dependencies} records the logical
order. The distinction between a fixed power saving and an estimate
uniform in the packet scale is essential throughout.
\begin{figure}[htbp]
\centering
\begin{tikzpicture}[
 node distance=9mm and 13mm,
 box/.style={draw=black!65,rounded corners=2pt,align=center,
             text width=39mm,minimum height=10mm,font=\small,inner sep=4pt},
 arrow/.style={-{Stealth[length=2mm]},semithick}]
 \node[box] (trans) {Transverse power saving};
 \node[box,below=of trans] (sparse) {Sparse packet estimate};
 \node[box,right=of trans] (frame) {Uniform bilinear frame estimate};
 \node[box,below=of frame] (flat) {Flat-frequency estimates};
 \node[box,right=of flat] (geom) {Rich-tube geometry};
 \node[box,below=14mm of flat] (dist) {Uniform dyadic\\distribution bound};
 \node[box,below=of dist] (sum) {Endpoint square summation};
 \node[box,below=of sum] (limit) {Gaussian limit at all times};
 \draw[arrow] (trans)--(sparse);
 \draw[arrow] (frame)--(flat);
 \draw[arrow] (frame.east) -- ++(6mm,0) |- ([yshift=2.5mm]dist.east);
 \draw[arrow] (sparse.south)|-(dist.west);
 \draw[arrow] (flat)--(dist);
 \draw[arrow] (geom.south)|-(dist.east);
 \draw[arrow] (dist)--(sum);
 \draw[arrow] (sum)--(limit);
\end{tikzpicture}
\caption{Dependencies of the proof. Small scale losses are allowed in the
transverse and sparse arguments; the frame and dyadic distribution estimates
must be uniform in the packet scale.}
\label{fig:dependencies}
\end{figure}

